% !PRINTED 140
\weberpage{140}

\weq{17}{K = \frac{1}{2}\int_0^1 \frac{d\xi}{\sqrt{\xi(1-\xi)(1-\varkappa^2\xi)}},}
\weq{18}{iK' = \frac{1}{2}\int_1^{\frac{1}{\varkappa^2}} \frac{d\xi}{\sqrt{\xi(1-\xi)(1-\varkappa^2\xi)}},}
\weq{19}{-K = \frac{1}{2}\int_{\frac{1}{\varkappa^2}}^{\infty} \frac{d\xi}{\sqrt{\xi(1-\xi)(1-\varkappa^2\xi)}},}
\weq{20}{-iK' = \frac{1}{2}\int_{\infty}^{0} \frac{d\xi}{\sqrt{\xi(1-\xi)(1-\varkappa^2\xi)}}.}
Die Werte von $\sqrt{\xi}$, $\sqrt{1-\xi}$, $\sqrt{1-\varkappa^2\xi}$ sind bei diesen Integrationen
 durch die Formeln (4) bestimmt. Wenn aber $\omega$ rein
imaginär und infolgedessen $\varkappa^2$ ein positiver echter Bruch und
$K$, $K'$ reell und positiv sind, so sind die Wege für $v$ der reellen
und imaginären Achse parallel, und mit Rücksicht auf die Bemerkungen
 am Schluß des \S\,25 zeigen die Formeln (4) folgendes:
\[
\begin{array}{lcccl}
\text{In (17) sind} & \sqrt{\xi}, & \sqrt{1-\xi}, & \sqrt{1-\varkappa^2\xi} & \text{reell und positiv,}\\
\text{„ (18) „} & \sqrt{\xi}, & i\sqrt{1-\xi}, & \sqrt{1-\varkappa^2\xi} & \text{„ „ „}\\
\text{„ (19) „} & \sqrt{\xi}, & i\sqrt{1-\xi}, & i\sqrt{1-\varkappa^2\xi} & \text{„ „ „}\\
\text{„ (20) „} & -i\sqrt{\xi}, & \sqrt{1-\xi}, & \sqrt{1-\varkappa^2\xi} & \text{„ „ „}
\end{array}
\]
und $\sqrt{\xi}$ hat auf keinem der Integrationswege ein Maximum oder
Minimum. Es können daher $K$, $K'$ durch die vier folgenden
\gesperrt{reellen Integrale} mit positiver Quadratwurzel erklärt werden:
\[
K = \frac{1}{2}\int_0^1 \frac{d\xi}{\sqrt{\xi(1-\xi)(1-\varkappa^2\xi)}} = \frac{1}{2}\int_{\frac{1}{\varkappa^2}}^{\infty} \frac{d\xi}{\sqrt{\xi(\xi-1)(\varkappa^2\xi-1)}},
\]
\[
K' = \frac{1}{2}\int_{-\infty}^{0} \frac{d\xi}{\sqrt{-\xi(1-\xi)(1-\varkappa^2\xi)}} = \frac{1}{2}\int_1^{\frac{1}{\varkappa^2}} \frac{d\xi}{\sqrt{\xi(\xi-1)(1-\varkappa^2\xi)}}.
\]
